% !TEX root =./ECE343_Practice_main.tex
%
% L7 practice set (Terminated Transmission Lines): the input impedance
% derivation, the special cases, stubs, the quarter-wave transformer, and power.

\fullwidth{\textbf{Learning Objectives 3.3 and 3.4: } I can use the transmission
line equation to calculate the voltage/current at an arbitrary point on a
specified transmission line. I can use the transmission line equation to match a
load to a source using a length of transmission line.
\\[4pt]\normalfont\small Show your work. Write a \textbf{Documentation} line at the
end (\textbf{None} if you did not collaborate). All lines are lossless. The load
is at $z = 0$ and $d$ is the distance from the load toward the source.}

\begin{questions}

%%%%%%%%%%%%%%%%%%%%%%%%%%%%%%%%%%%%%%%%%
\question \textbf{Deriving the input impedance.} The phasor voltage and current
on a line terminated in $\myz{L}$ are
\eq{
	\tilde{V}(z) &= V_0^+\lp e^{-j\beta z} + \Gamma e^{j\beta z}\rp, &
	\tilde{I}(z) &= \dfrac{V_0^+}{\myz{0}}\lp e^{-j\beta z} - \Gamma e^{j\beta z}\rp
}
\begin{parts}

	\part Write $\myz{in}(d) = \tilde{V}(-d)/\tilde{I}(-d)$, substitute
	$\Gamma = \frac{\myz{L} - \myz{0}}{\myz{L} + \myz{0}}$, and show that
	\[ \myz{in}(d) = \myz{0}\,\dfrac{\myz{L} + j\myz{0}\tan\beta d}{\myz{0} + j\myz{L}\tan\beta d}. \]
	\begin{solution}[3.2in]
	\eq{
		\myz{in} &= \myz{0}\,\dfrac{e^{j\beta d} + \Gamma e^{-j\beta d}}{e^{j\beta d} - \Gamma e^{-j\beta d}}
	}
	Multiply the numerator and denominator by $\myz{L} + \myz{0}$:
	\eq{
		\myz{in} &= \myz{0}\,\dfrac{(\myz{L} + \myz{0})e^{j\beta d} + (\myz{L} - \myz{0})e^{-j\beta d}}{(\myz{L} + \myz{0})e^{j\beta d} - (\myz{L} - \myz{0})e^{-j\beta d}} \\
		&= \myz{0}\,\dfrac{\myz{L}\lp e^{j\beta d} + e^{-j\beta d}\rp + \myz{0}\lp e^{j\beta d} - e^{-j\beta d}\rp}{\myz{0}\lp e^{j\beta d} + e^{-j\beta d}\rp + \myz{L}\lp e^{j\beta d} - e^{-j\beta d}\rp} \\
		&= \myz{0}\,\dfrac{2\myz{L}\cos\beta d + j2\myz{0}\sin\beta d}{2\myz{0}\cos\beta d + j2\myz{L}\sin\beta d}
	}
	using Euler's identity. Dividing the numerator and denominator by
	$2\cos\beta d$ gives the result.
	\end{solution}

	\part Set $\myz{L} = 0$ and then $\myz{L} \to \infty$ to find the input
	impedance of a shorted line and of an open line.
	\begin{solution}[1.2in]
	\eq{
		\myz{L} = 0: \quad \myz{in} &= \myz{0}\,\dfrac{j\myz{0}\tan\beta d}{\myz{0}} = j\myz{0}\tan\beta d \\
		\myz{L} \to \infty: \quad \myz{in} &= \myz{0}\,\dfrac{\myz{L}}{j\myz{L}\tan\beta d} = -j\myz{0}\cot\beta d
	}
	For the open line we divide the numerator and denominator by $\myz{L}$ before
	taking the limit.
	\end{solution}

\end{parts}

\newpage
%%%%%%%%%%%%%%%%%%%%%%%%%%%%%%%%%%%%%%%%%
\question \textbf{One load, three lengths.} A $100\ \ohms$ resistor terminates
a $50\ \ohms$ line.
\begin{parts}

	\begin{minipage}[t]{0.58\linewidth}
		\part Find $\myz{in}$ for a line $\lambda/8$ long.
		\begin{solution}[1.6in]
		$\beta d = 45^\circ$, so $\tan\beta d = 1$:
		\eq{
			\myz{in} &= 50\,\dfrac{100 + j50}{50 + j100} = 50\,\dfrac{(100 + j50)(50 - j100)}{50^2 + 100^2} \\
			&= 50\,\dfrac{10{,}000 - j7{,}500}{12{,}500} = 40 - j30\ \ohms
		}
		\end{solution}
	\end{minipage}\hfill%
	\begin{minipage}[t]{0.37\linewidth}
		\raggedleft \vspace{12pt}
		$\myz{in}$ = \ansbox[1.3in]{$40 - j30\ \ohms$}
	\end{minipage}

	\begin{minipage}[t]{0.58\linewidth}
		\part Find $\myz{in}$ for lines $\lambda/4$ and $\lambda/2$ long.
		\begin{solution}[1in]
		\eq{
			\lambda/4: \quad \myz{in} &= \dfrac{\myz{0}^2}{\myz{L}} = \dfrac{2500}{100} = 25\ \ohms \\
			\lambda/2: \quad \myz{in} &= \myz{L} = 100\ \ohms
		}
		\end{solution}
	\end{minipage}\hfill%
	\begin{minipage}[t]{0.37\linewidth}
		\raggedleft \vspace{12pt}
		$\lambda/4$: \ansbox[1.3in]{$25\ \ohms$}

		\vspace{4pt}
		$\lambda/2$: \ansbox[1.3in]{$100\ \ohms$}
	\end{minipage}

	\part The resistor absorbs real power, but the input impedance in part (a)
	has a reactive part. Explain in one or two sentences how both can be true.
	\begin{solution}[1in]
	The reactance in $\myz{in}$ comes from the standing wave on the line, which
	stores energy that flows back and forth between the line's inductance and
	capacitance. The lossless line dissipates nothing, so the real power that
	enters at the input is the real power the resistor absorbs.
	\end{solution}

\end{parts}

\newpage
%%%%%%%%%%%%%%%%%%%%%%%%%%%%%%%%%%%%%%%%%
\question \textbf{Stubs.} A circuit at $2\ \ghz$ needs a $5\ \text{nH}$ inductor
and a $1\ \text{pF}$ capacitor. Build each from $50\ \ohms$ line in which waves
travel at $2\times10^{8}$ m/s.
\begin{parts}

	\begin{minipage}[t]{0.58\linewidth}
		\part Find the wavelength in the line.
		\begin{solution}[0.7in]
		\eq{ \lambda = \dfrac{2\times10^{8}}{2\times10^{9}} = 0.1\ \m = 10\ \text{cm} }
		\end{solution}
	\end{minipage}\hfill%
	\begin{minipage}[t]{0.37\linewidth}
		\raggedleft \vspace{12pt}
		$\lambda$ = \ansbox[1.3in]{$10\ \text{cm}$}
	\end{minipage}

	\begin{minipage}[t]{0.58\linewidth}
		\part Find the length of the shortest shorted stub that acts as the
		inductor.
		\begin{solution}[1.3in]
		\eq{
			\omega L &= 2\pi \times 2\times10^{9} \times 5\times10^{-9} = 62.8\ \ohms \\
			\tan\beta\ell &= \dfrac{62.8}{50} = 1.257 \quad\Rightarrow\quad \beta\ell = 51.5^\circ \\
			\ell &= \dfrac{51.5^\circ}{360^\circ}\lambda = 0.143\lambda = 1.43\ \text{cm}
		}
		\end{solution}
	\end{minipage}\hfill%
	\begin{minipage}[t]{0.37\linewidth}
		\raggedleft \vspace{12pt}
		$\ell$ = \ansbox[1.3in]{$1.43\ \text{cm}$}
	\end{minipage}

	\begin{minipage}[t]{0.58\linewidth}
		\part Find the length of the shortest open stub that acts as the
		capacitor.
		\begin{solution}[1.3in]
		\eq{
			\dfrac{1}{\omega C} &= \dfrac{1}{2\pi \times 2\times10^{9} \times 10^{-12}} = 79.6\ \ohms \\
			\cot\beta\ell &= \dfrac{79.6}{50} = 1.592 \quad\Rightarrow\quad \beta\ell = 32.1^\circ \\
			\ell &= \dfrac{32.1^\circ}{360^\circ}\lambda = 0.089\lambda = 0.89\ \text{cm}
		}
		\end{solution}
	\end{minipage}\hfill%
	\begin{minipage}[t]{0.37\linewidth}
		\raggedleft \vspace{12pt}
		$\ell$ = \ansbox[1.3in]{$0.89\ \text{cm}$}
	\end{minipage}

	\part Your stub in part (b) is accidentally cut $2.5\ \text{cm}$ longer
	than designed. What does it act like now? Explain without computing.
	\begin{solution}[0.9in]
	$2.5\ \text{cm}$ is $\lambda/4$, and a shorted line between $\lambda/4$ and
	$\lambda/2$ long has $\tan\beta\ell < 0$: it acts as a capacitor.
	\end{solution}

\end{parts}

\newpage
%%%%%%%%%%%%%%%%%%%%%%%%%%%%%%%%%%%%%%%%%
\question \textbf{Quarter-wave transformer.} A $200\ \ohms$ antenna is fed at
$500\ \mhz$ from a $50\ \ohms$ line. Waves travel at $2\times10^{8}$ m/s in
every line.
\begin{parts}

	\begin{minipage}[t]{0.58\linewidth}
		\part Find the characteristic impedance and the length of the
		quarter-wave section that matches the antenna.
		\begin{solution}[1.2in]
		\eq{
			\myz{1} &= \sqrt{\myz{0}\,R_L} = \sqrt{50 \times 200} = 100\ \ohms \\
			\lambda &= \dfrac{2\times10^{8}}{5\times10^{8}} = 40\ \text{cm}, \quad \dfrac{\lambda}{4} = 10\ \text{cm}
		}
		\end{solution}
	\end{minipage}\hfill%
	\begin{minipage}[t]{0.37\linewidth}
		\raggedleft \vspace{12pt}
		$\myz{1}$ = \ansbox[1.3in]{$100\ \ohms$}

		\vspace{4pt}
		length = \ansbox[1.3in]{$10\ \text{cm}$}
	\end{minipage}

	\begin{minipage}[t]{0.58\linewidth}
		\part The same network is used at $1\ \ghz$. Find the impedance the
		$50\ \ohms$ line sees and the reflection coefficient there.
		\begin{solution}[1.2in]
		At $1\ \ghz$ the $10\ \text{cm}$ section is $\lambda/2$ long, so it
		repeats the load:
		\eq{
			\myz{in} &= 200\ \ohms \\
			\Gamma &= \dfrac{200 - 50}{200 + 50} = 0.6
		}
		\end{solution}
	\end{minipage}\hfill%
	\begin{minipage}[t]{0.37\linewidth}
		\raggedleft \vspace{12pt}
		$\myz{in}$ = \ansbox[1.3in]{$200\ \ohms$}

		\vspace{4pt}
		$\Gamma$ = \ansbox[1.3in]{$0.6$}
	\end{minipage}

\end{parts}

\newpage
%%%%%%%%%%%%%%%%%%%%%%%%%%%%%%%%%%%%%%%%%
\question \textbf{Power delivered.} An incident wave of amplitude
$\vert V_0^+\vert = 10\ \text{V}$ travels on a $50\ \ohms$ line toward a
$100\ \ohms$ load.
\begin{parts}

	\begin{minipage}[t]{0.58\linewidth}
		\part Find the incident power.
		\begin{solution}[0.8in]
		\eq{ P_\text{inc} = \dfrac{\vert V_0^+\vert^2}{2\myz{0}} = \dfrac{100}{100} = 1\ \text{W} }
		\end{solution}
	\end{minipage}\hfill%
	\begin{minipage}[t]{0.37\linewidth}
		\raggedleft \vspace{12pt}
		$P_\text{inc}$ = \ansbox[1.3in]{$1\ \text{W}$}
	\end{minipage}

	\begin{minipage}[t]{0.58\linewidth}
		\part Find the power the load absorbs.
		\begin{solution}[1in]
		\eq{
			\Gamma &= \dfrac{100 - 50}{100 + 50} = \dfrac{1}{3} \\
			P_L &= P_\text{inc}\lp 1 - \vert\Gamma\vert^2\rp = 1 \times \dfrac{8}{9} = 0.889\ \text{W}
		}
		\end{solution}
	\end{minipage}\hfill%
	\begin{minipage}[t]{0.37\linewidth}
		\raggedleft \vspace{12pt}
		$P_L$ = \ansbox[1.3in]{$0.889\ \text{W}$}
	\end{minipage}

	\begin{minipage}[t]{0.58\linewidth}
		\part Check part (b) directly: find the load voltage
		$\tilde{V}(0) = V_0^+(1 + \Gamma)$ and compute
		$\frac{1}{2}\vert\tilde{V}(0)\vert^2/R_L$.
		\begin{solution}[1in]
		\eq{
			\vert\tilde{V}(0)\vert &= 10 \times \dfrac{4}{3} = 13.3\ \text{V} \\
			P_L &= \dfrac{13.3^2}{2 \times 100} = 0.889\ \text{W}
		}
		\end{solution}
	\end{minipage}\hfill%
	\begin{minipage}[t]{0.37\linewidth}
		\raggedleft \vspace{12pt}
		$P_L$ = \ansbox[1.3in]{$0.889\ \text{W}$}
	\end{minipage}

\end{parts}

\vspace{1.5em}
\noindent\textbf{Documentation:} \rule[-2pt]{0.6\linewidth}{0.4pt}

\end{questions}
